\documentclass[aspectratio=169,10pt]{beamer}
\makeatletter\def\input@path{{../../estilo/}{../estilo/}{estilo/}}\makeatother
\usepackage{devmobile}
\modulo[Flutter]{15}{Publicação: Google Play e App Store}{Identidade do app, builds de release, assinatura, Play Console, App Store Connect, TestFlight, revisão, automação com fastlane e operação depois do lançamento}

\begin{document}
\capa

\begin{frame}{Objetivos do módulo}
  \begin{conceito}[Ao final deste módulo você será capaz de]
    \begin{itemize}
      \item Explicar o caminho do código até o aparelho do usuário e o papel da assinatura digital.
      \item Configurar nome, ícone, \emph{splash}, versão e ambientes de um app Flutter.
      \item Gerar builds de release ofuscados para Android e iOS.
      \item Publicar na Google Play passando pelas faixas de teste.
      \item Publicar na App Store passando pelo TestFlight e pela revisão.
      \item Automatizar a entrega com fastlane e operar o app depois do lançamento.
    \end{itemize}
  \end{conceito}
\end{frame}

\sumario

% =============================================================================
%  Fundamentos compartilhados --- Módulo 15: distribuição de apps
%  Incluído pelas trilhas Kotlin e Flutter com % =============================================================================
%  Fundamentos compartilhados --- Módulo 15: distribuição de apps
%  Incluído pelas trilhas Kotlin e Flutter com % =============================================================================
%  Fundamentos compartilhados --- Módulo 15: distribuição de apps
%  Incluído pelas trilhas Kotlin e Flutter com \input{../fundamentos/f15-publicacao}
% =============================================================================
\section{Fundamentos da distribuição de apps}

\begin{frame}{Do código ao aparelho do usuário}
  \centering
  \begin{tikzpicture}[node distance=2.4mm,every node/.style={font=\scriptsize,minimum height=9mm,minimum width=1.62cm}]
    \node[caixa] (c) {Código\\ \tiny versionado};
    \node[caixapetroleo,right=of c] (b) {Build de\\ release};
    \node[caixapetroleo,right=of b] (a) {Assinatura};
    \node[caixaambar,right=of a] (e) {Envio\\ à loja};
    \node[caixaambar,right=of e] (r) {Revisão\\ e políticas};
    \node[caixa,right=of r] (t) {Faixas\\ de teste};
    \node[caixaescura,right=of t] (p) {Produção};
    \foreach \x/\y in {c/b,b/a,a/e,e/r,r/t,t/p} \draw[seta] (\x) -- (\y);
    \draw[setaambar,rounded corners] (p.south) -- ++(0,-6mm) -| node[rotulo,pos=0.25,below]{\tiny monitorar falhas e avaliações $\rightarrow$ nova versão} (c.south);
  \end{tikzpicture}
  \vspace{6pt}
  \begin{columns}[T]
    \begin{column}{0.48\textwidth}
      \scriptsize
      \begin{tabularx}{\linewidth}{@{}l L L@{}}
        \toprule
        & \cabecalho{Debug} & \cabecalho{Release} \\
        \midrule
        Objetivo & desenvolver & distribuir \\
        Código & com verificações e depuração & otimizado, reduzido e ofuscável \\
        Desempenho & mais lento & o real \\
        Assinatura & chave de depuração automática & sua chave de release \\
        \bottomrule
      \end{tabularx}
    \end{column}
    \begin{column}{0.48\textwidth}
      \begin{definicao}[Build de release]
        Versão do app compilada para o usuário final: sem ferramentas de depuração, com código otimizado e assinada com a chave do desenvolvedor.
      \end{definicao}
      \begin{atencao}
        Teste sempre o build de release antes de publicar: ofuscação e otimizações podem expor defeitos que não aparecem em debug.
      \end{atencao}
    \end{column}
  \end{columns}
\end{frame}

\begin{frame}{Assinatura digital: a teoria}
  \begin{columns}[c]
    \begin{column}{0.54\textwidth}
      \centering
      \begin{tikzpicture}[node distance=4mm,every node/.style={font=\scriptsize}]
        \node[caixa,minimum width=1.6cm] (app) {App};
        \node[caixa,right=of app,minimum width=1.6cm] (h) {Hash\\ \tiny resumo único};
        \node[caixaescura,right=of h,minimum width=2cm] (s) {Assinatura};
        \node[caixaambar,above=5mm of s,minimum width=2cm] (pk) {Chave \textbf{privada}\\ \tiny só o dev tem};
        \draw[seta] (app) -- (h); \draw[seta] (h) -- (s); \draw[seta] (pk) -- (s);
        \node[caixa,below=9mm of app,minimum width=1.6cm] (app2) {App\\ \tiny recebido};
        \node[caixa,right=of app2,minimum width=1.6cm] (h2) {Hash\\ \tiny recalculado};
        \node[caixapetroleo,right=of h2,minimum width=2cm] (v) {Confere?};
        \node[caixa,below=5mm of v,minimum width=2cm] (pub) {Chave \textbf{pública}\\ \tiny no certificado};
        \draw[seta] (app2) -- (h2); \draw[seta] (h2) -- (v); \draw[seta] (pub) -- (v);
        \draw[setaambar] (s.south) -- node[rotulo,right]{\tiny vai junto} (v.north);
        \node[rotulo,left=1mm of app,font=\tiny,rotate=90,anchor=south] {quem publica};
        \node[rotulo,left=1mm of app2,font=\tiny,rotate=90,anchor=south] {sistema};
      \end{tikzpicture}
    \end{column}
    \begin{column}{0.44\textwidth}
      \begin{definicao}[Assinatura digital]
        Resumo (\emph{hash}) do pacote cifrado com a \textbf{chave privada} do desenvolvedor. Qualquer um verifica com a \textbf{chave pública}, mas só quem tem a privada consegue assinar.
      \end{definicao}
      \begin{conceito}[O que ela garante]
        \textbf{Integridade}: o pacote não foi alterado.\\
        \textbf{Autenticidade}: a atualização vem do mesmo autor do app instalado.
      \end{conceito}
      \begin{analogia}
        O lacre de cera com o sinete do remetente: qualquer um reconhece o selo, ninguém o falsifica sem o anel.
      \end{analogia}
    \end{column}
  \end{columns}
\end{frame}

\begin{frame}{Assinatura em cada plataforma}
  \begin{columns}[T]
    \begin{column}{0.5\textwidth}
      \begin{conceito}[Android]
        O pacote é assinado com uma chave guardada num \emph{keystore}. Com o \textbf{Play App Signing}, você assina com a \textbf{chave de upload}; o Google confere e reassina com a \textbf{chave do app}, que fica guardada com ele.
      \end{conceito}
      \centering
      \begin{tikzpicture}[node distance=9mm,every node/.style={font=\tiny,align=center}]
        \node[caixa,text width=1.5cm] (d) {Dev\\ chave de upload};
        \node[caixapetroleo,text width=1.5cm,right=of d] (g) {Google Play\\ chave do app};
        \node[caixa,right=of g] (u) {Aparelho};
        \draw[seta] (d) -- node[rotulo,above]{AAB} (g);
        \draw[seta] (g) -- node[rotulo,above]{APKs} (u);
      \end{tikzpicture}
      \begin{dica}
        Perdeu a chave de upload? Com o Play App Signing, o Google permite trocá-la. Sem ele, perder a chave impede atualizar o app.
      \end{dica}
    \end{column}
    \begin{column}{0.48\textwidth}
      \begin{conceito}[iOS]
        Todo app precisa ser assinado para rodar num iPhone. Três peças: o \textbf{certificado} (sua identidade), o \textbf{App ID} e o \textbf{perfil de provisionamento}, que liga certificado, app, capacidades e aparelhos.
      \end{conceito}
      \begin{dica}
        O Xcode gerencia tudo isso com \emph{Automatically manage signing}.
      \end{dica}
      \begin{atencao}
        Chaves e certificados são segredos: nunca no Git; guarde cópias num cofre de senhas.
      \end{atencao}
    \end{column}
  \end{columns}
\end{frame}

\begin{frame}[fragile]{Versionamento}
  \begin{columns}[T]
    \begin{column}{0.5\textwidth}
      \begin{definicao}[Versionamento semântico (SemVer)]
        \textbf{MAJOR.MINOR.PATCH}: \textbf{MAJOR} muda com quebras de compatibilidade, \textbf{MINOR} com funcionalidades novas compatíveis e \textbf{PATCH} com correções.
      \end{definicao}
      \centering
      \begin{tikzpicture}[every node/.style={font=\scriptsize}]
        \node[caixapetroleo] (m) {2};
        \node[right=0.5mm of m] (p1) {.};
        \node[caixaambar,right=0.5mm of p1] (n) {4};
        \node[right=0.5mm of n] (p2) {.};
        \node[caixa,right=0.5mm of p2] (c) {1};
        \node[right=0.5mm of c] (p3) {+};
        \node[caixaescura,right=0.5mm of p3] (b) {27};
        \node[rotulo,below=1mm of m,font=\tiny] {major};
        \node[rotulo,below=1mm of n,font=\tiny] {minor};
        \node[rotulo,below=1mm of c,font=\tiny] {patch};
        \node[rotulo,below=1mm of b,font=\tiny,align=center] {número\\do build};
      \end{tikzpicture}
    \end{column}
    \begin{column}{0.48\textwidth}
      \scriptsize
      \begin{tabularx}{\linewidth}{@{}l L L@{}}
        \toprule
        & \cabecalho{Para pessoas} & \cabecalho{Para a loja} \\
        \midrule
        Android & \cod{versionName} & \cod{versionCode} \\
        iOS & \cod{CFBundleShort\-VersionString} & \cod{CFBundleVersion} \\
        Flutter & \cod{version: 2.4.1+27} & (o número após o \cod{+}) \\
        \bottomrule
      \end{tabularx}
      \begin{atencao}[Regra das lojas]
        O número do build \textbf{sempre aumenta}. Não existe ``voltar versão'': para desfazer um problema, publica-se uma versão nova com número maior.
      \end{atencao}
    \end{column}
  \end{columns}
\end{frame}

\begin{frame}{Formatos de pacote e tamanho do app}
  \begin{columns}[T]
    \begin{column}{0.52\textwidth}
      \scriptsize
      \begin{tabularx}{\linewidth}{@{}l L@{}}
        \toprule
        \cabecalho{Formato} & \cabecalho{O que é} \\
        \midrule
        \textbf{APK} & pacote Android instalável, com tudo para todos os aparelhos \\
        \textbf{AAB} & \emph{Android App Bundle}: formato de \textbf{envio} à Play, que gera APKs sob medida para cada aparelho \\
        \textbf{IPA} & pacote iOS enviado à App Store; a Apple também entrega variantes por aparelho (\emph{App Thinning}) \\
        \bottomrule
      \end{tabularx}
      \centering
      \begin{tikzpicture}[node distance=2mm,every node/.style={font=\tiny}]
        \node[caixapetroleo,minimum width=1.6cm] (aab) {AAB};
        \node[caixa,right=8mm of aab,yshift=5mm] (a1) {arm64 + xxhdpi + pt};
        \node[caixa,right=8mm of aab,yshift=-5mm] (a2) {armeabi + hdpi + en};
        \draw[seta] (aab) -- (a1.west); \draw[seta] (aab) -- (a2.west);
      \end{tikzpicture}
    \end{column}
    \begin{column}{0.46\textwidth}
      \begin{conceito}[Por que o tamanho importa]
        Apps grandes demoram a baixar, ocupam espaço e perdem instalações, principalmente em planos de dados limitados.
      \end{conceito}
      \begin{pordentro}[Como os builds encolhem]
        \textbf{\emph{Tree shaking}} e \textbf{\emph{shrinking}} removem código não usado; a \textbf{ofuscação} troca nomes por símbolos curtos; recursos são comprimidos e separados por aparelho.
      \end{pordentro}
      \begin{atencao}
        Código ofuscado gera relatórios de falha ilegíveis. Guarde os \textbf{símbolos de depuração} de cada versão para decodificá-los.
      \end{atencao}
    \end{column}
  \end{columns}
\end{frame}

\begin{frame}{As duas grandes lojas}
  \scriptsize
  \begin{tabularx}{\textwidth}{@{}p{2.6cm} L L@{}}
    \toprule
    & \cabecalho{Google Play} & \cabecalho{App Store} \\
    \midrule
    Conta & taxa única de US\$\,25; pessoal ou organização & Apple Developer Program, US\$\,99 por ano (isenção para instituições de ensino e sem fins lucrativos) \\
    Revisão & automatizada e humana; de horas a dias (contas novas esperam mais) & humana (\emph{App Review}); em geral de um a dois dias \\
    Testes & faixas interna, fechada e aberta & TestFlight interno (até 100) e externo (até 10\,000) \\
    Lançamento gradual & percentual definido por você & \emph{phased release} automático em 7 dias \\
    Requisito anual & \cod{targetSdk} 36 desde 31/08/2026 & Xcode 26 e SDK do iOS 26 desde 28/04/2026 \\
    Comissão & de 15\,\% a 30\,\% sobre vendas digitais, conforme receita, produto e país & idem \\
    \bottomrule
  \end{tabularx}
  \vspace{4pt}
  \begin{dica}
    Regras de comissão e de pagamento externo mudam com leis e decisões judiciais (União Europeia, EUA, Japão, Coreia). Confira antes de definir o modelo de negócio (Módulo 20).
  \end{dica}
\end{frame}

\begin{frame}{Privacidade desde a concepção e LGPD}
  \begin{columns}[T]
    \begin{column}{0.5\textwidth}
      \begin{definicao}[LGPD (Lei 13.709/2018)]
        Regula o tratamento de \textbf{dados pessoais} no Brasil: tudo que identifica ou pode identificar uma pessoa (nome, e-mail, matrícula, localização, identificador do aparelho).
      \end{definicao}
      \begin{conceito}[Princípios que mais afetam um app]
        \textbf{Finalidade} (para quê), \textbf{necessidade} (o mínimo), \textbf{transparência} (explicar), \textbf{segurança} (proteger) e uma \textbf{base legal} para cada uso (consentimento, execução de contrato\ldots).
      \end{conceito}
    \end{column}
    \begin{column}{0.48\textwidth}
      \begin{conceito}[O que as lojas exigem]
        \begin{itemize}
          \item Política de privacidade publicada.
          \item Declaração de dados coletados: \emph{Data safety} (Play) e \emph{App Privacy} (App Store).
          \item Textos de uso para cada permissão.
          \item \textbf{Exclusão de conta} dentro do app, se ele permite criar conta.
        \end{itemize}
      \end{conceito}
      \begin{atencao}
        A declaração deve incluir o que os SDKs de terceiros coletam (analytics, anúncios, Crashlytics). Não é assessoria jurídica: consulte o encarregado de dados (DPO) da instituição.
      \end{atencao}
    \end{column}
  \end{columns}
\end{frame}

\begin{frame}{Gestão de lançamentos}
  \begin{columns}[c]
    \begin{column}{0.52\textwidth}
      \centering
      \begin{tikzpicture}[x=1.0cm,y=0.45cm]
        \draw[cinzamedio,-{Stealth}] (0,0) -- (6.3,0) node[right,font=\tiny]{dias};
        \draw[cinzamedio,-{Stealth}] (0,0) -- (0,5.6) node[above,font=\tiny]{\% dos usuários};
        \draw[azulpetroleo,line width=1.4pt] (0,0.3) -- (1,0.3) -- (1,1) -- (2,1) -- (2,2.5) -- (3.5,2.5) -- (3.5,5) -- (6,5);
        \node[font=\tiny,azulpetroleo] at (0.5,0.9) {1\%};
        \node[font=\tiny,azulpetroleo] at (1.5,1.6) {5\%};
        \node[font=\tiny,azulpetroleo] at (2.75,3.1) {20\%};
        \node[font=\tiny,azulpetroleo] at (4.75,5.6) {100\%};
        \foreach \x in {1,2,3.5} \fill[ambar] (\x,0) circle (2.2pt);
        \node[rotulo,font=\tiny,align=center] at (2.2,-1.6) {\textcolor{ambar}{$\bullet$} portões: falhas e avaliações\\ dentro do limite? então avança};
      \end{tikzpicture}
    \end{column}
    \begin{column}{0.46\textwidth}
      \begin{definicao}[Lançamento gradual]
        Liberar a versão para uma fração crescente dos usuários, avançando só se os indicadores de saúde continuarem bons.
      \end{definicao}
      \begin{conceito}[Ferramentas do processo]
        \begin{itemize}
          \item \textbf{Trem de versões}: lançar em datas fixas (a cada duas semanas, por exemplo).
          \item \textbf{\emph{Feature flags}}: ligar recursos remotamente, sem nova versão.
          \item \textbf{Versão mínima}: obrigar a atualização quando a antiga é insegura.
          \item \textbf{Notas de versão} claras para usuários e revisores.
        \end{itemize}
      \end{conceito}
    \end{column}
  \end{columns}
\end{frame}

\begin{frame}{ASO: o funil da loja}
  \begin{columns}[c]
    \begin{column}{0.46\textwidth}
      \centering
      \begin{tikzpicture}[every node/.style={font=\scriptsize\bfseries,text=white}]
        \fill[azulpetroleo] (-2.6,3.2) -- (2.6,3.2) -- (2.2,2.5) -- (-2.2,2.5) -- cycle;
        \node at (0,2.85) {Impressões na busca};
        \fill[azulprofundo] (-2.15,2.4) -- (2.15,2.4) -- (1.75,1.7) -- (-1.75,1.7) -- cycle;
        \node at (0,2.05) {Visitas à ficha};
        \fill[ambar!90!black] (-1.7,1.6) -- (1.7,1.6) -- (1.3,0.9) -- (-1.3,0.9) -- cycle;
        \node at (0,1.25) {Instalações};
        \fill[ambar] (-1.25,0.8) -- (1.25,0.8) -- (0.85,0.1) -- (-0.85,0.1) -- cycle;
        \node[text=azulpetroleo] at (0,0.45) {Uso ativo};
      \end{tikzpicture}
    \end{column}
    \begin{column}{0.52\textwidth}
      \begin{definicao}[ASO --- \emph{App Store Optimization}]
        Conjunto de práticas para o app ser \textbf{encontrado} na busca da loja e \textbf{escolhido} por quem vê a ficha.
      \end{definicao}
      \scriptsize
      \begin{tabularx}{\linewidth}{@{}l L@{}}
        \toprule
        \cabecalho{Etapa} & \cabecalho{O que influencia} \\
        \midrule
        Ser encontrado & título, subtítulo ou descrição curta, palavras-chave, categoria \\
        Ser escolhido & ícone, capturas de tela, vídeo, nota e avaliações \\
        Ser mantido & qualidade, desempenho, respostas às avaliações \\
        \bottomrule
      \end{tabularx}
      \begin{dica}
        As duas lojas permitem testes A/B da ficha (ícone, capturas, textos).
      \end{dica}
    \end{column}
  \end{columns}
\end{frame}

\begin{frame}{Depois do lançamento: saúde do app}
  \begin{columns}[T]
    \begin{column}{0.5\textwidth}
      \scriptsize
      \begin{tabularx}{\linewidth}{@{}p{2.3cm} L@{}}
        \toprule
        \cabecalho{Indicador} & \cabecalho{O que mede} \\
        \midrule
        Usuários sem falhas & \% de usuários que não tiveram nenhum \emph{crash} no período \\
        Taxa de ANR & \% de sessões em que a tela travou (Android) \\
        Tempo de inicialização & do toque no ícone à primeira tela útil \\
        Retenção D1/D7/D30 & \% de usuários que voltam 1, 7 e 30 dias após instalar \\
        Nota e avaliações & percepção dos usuários; comentários apontam defeitos \\
        \bottomrule
      \end{tabularx}
    \end{column}
    \begin{column}{0.48\textwidth}
      \begin{atencao}[Limites do Android vitals]
        Acima de 1,09\,\% de falhas percebidas ou 0,47\,\% de ANRs percebidos, o app é considerado de ``mau comportamento'' e perde visibilidade na Play.
      \end{atencao}
      \begin{conceito}[Ciclo contínuo]
        Monitorar $\rightarrow$ priorizar $\rightarrow$ corrigir $\rightarrow$ lançar gradualmente $\rightarrow$ monitorar. Falhas e avaliações alimentam o \emph{backlog} da próxima versão.
      \end{conceito}
      \begin{dica}
        Responda às avaliações com educação e objetividade: a resposta é pública e mostra que o app é cuidado.
      \end{dica}
    \end{column}
  \end{columns}
\end{frame}

% =============================================================================
\section{Fundamentos da distribuição de apps}

\begin{frame}{Do código ao aparelho do usuário}
  \centering
  \begin{tikzpicture}[node distance=2.4mm,every node/.style={font=\scriptsize,minimum height=9mm,minimum width=1.62cm}]
    \node[caixa] (c) {Código\\ \tiny versionado};
    \node[caixapetroleo,right=of c] (b) {Build de\\ release};
    \node[caixapetroleo,right=of b] (a) {Assinatura};
    \node[caixaambar,right=of a] (e) {Envio\\ à loja};
    \node[caixaambar,right=of e] (r) {Revisão\\ e políticas};
    \node[caixa,right=of r] (t) {Faixas\\ de teste};
    \node[caixaescura,right=of t] (p) {Produção};
    \foreach \x/\y in {c/b,b/a,a/e,e/r,r/t,t/p} \draw[seta] (\x) -- (\y);
    \draw[setaambar,rounded corners] (p.south) -- ++(0,-6mm) -| node[rotulo,pos=0.25,below]{\tiny monitorar falhas e avaliações $\rightarrow$ nova versão} (c.south);
  \end{tikzpicture}
  \vspace{6pt}
  \begin{columns}[T]
    \begin{column}{0.48\textwidth}
      \scriptsize
      \begin{tabularx}{\linewidth}{@{}l L L@{}}
        \toprule
        & \cabecalho{Debug} & \cabecalho{Release} \\
        \midrule
        Objetivo & desenvolver & distribuir \\
        Código & com verificações e depuração & otimizado, reduzido e ofuscável \\
        Desempenho & mais lento & o real \\
        Assinatura & chave de depuração automática & sua chave de release \\
        \bottomrule
      \end{tabularx}
    \end{column}
    \begin{column}{0.48\textwidth}
      \begin{definicao}[Build de release]
        Versão do app compilada para o usuário final: sem ferramentas de depuração, com código otimizado e assinada com a chave do desenvolvedor.
      \end{definicao}
      \begin{atencao}
        Teste sempre o build de release antes de publicar: ofuscação e otimizações podem expor defeitos que não aparecem em debug.
      \end{atencao}
    \end{column}
  \end{columns}
\end{frame}

\begin{frame}{Assinatura digital: a teoria}
  \begin{columns}[c]
    \begin{column}{0.54\textwidth}
      \centering
      \begin{tikzpicture}[node distance=4mm,every node/.style={font=\scriptsize}]
        \node[caixa,minimum width=1.6cm] (app) {App};
        \node[caixa,right=of app,minimum width=1.6cm] (h) {Hash\\ \tiny resumo único};
        \node[caixaescura,right=of h,minimum width=2cm] (s) {Assinatura};
        \node[caixaambar,above=5mm of s,minimum width=2cm] (pk) {Chave \textbf{privada}\\ \tiny só o dev tem};
        \draw[seta] (app) -- (h); \draw[seta] (h) -- (s); \draw[seta] (pk) -- (s);
        \node[caixa,below=9mm of app,minimum width=1.6cm] (app2) {App\\ \tiny recebido};
        \node[caixa,right=of app2,minimum width=1.6cm] (h2) {Hash\\ \tiny recalculado};
        \node[caixapetroleo,right=of h2,minimum width=2cm] (v) {Confere?};
        \node[caixa,below=5mm of v,minimum width=2cm] (pub) {Chave \textbf{pública}\\ \tiny no certificado};
        \draw[seta] (app2) -- (h2); \draw[seta] (h2) -- (v); \draw[seta] (pub) -- (v);
        \draw[setaambar] (s.south) -- node[rotulo,right]{\tiny vai junto} (v.north);
        \node[rotulo,left=1mm of app,font=\tiny,rotate=90,anchor=south] {quem publica};
        \node[rotulo,left=1mm of app2,font=\tiny,rotate=90,anchor=south] {sistema};
      \end{tikzpicture}
    \end{column}
    \begin{column}{0.44\textwidth}
      \begin{definicao}[Assinatura digital]
        Resumo (\emph{hash}) do pacote cifrado com a \textbf{chave privada} do desenvolvedor. Qualquer um verifica com a \textbf{chave pública}, mas só quem tem a privada consegue assinar.
      \end{definicao}
      \begin{conceito}[O que ela garante]
        \textbf{Integridade}: o pacote não foi alterado.\\
        \textbf{Autenticidade}: a atualização vem do mesmo autor do app instalado.
      \end{conceito}
      \begin{analogia}
        O lacre de cera com o sinete do remetente: qualquer um reconhece o selo, ninguém o falsifica sem o anel.
      \end{analogia}
    \end{column}
  \end{columns}
\end{frame}

\begin{frame}{Assinatura em cada plataforma}
  \begin{columns}[T]
    \begin{column}{0.5\textwidth}
      \begin{conceito}[Android]
        O pacote é assinado com uma chave guardada num \emph{keystore}. Com o \textbf{Play App Signing}, você assina com a \textbf{chave de upload}; o Google confere e reassina com a \textbf{chave do app}, que fica guardada com ele.
      \end{conceito}
      \centering
      \begin{tikzpicture}[node distance=9mm,every node/.style={font=\tiny,align=center}]
        \node[caixa,text width=1.5cm] (d) {Dev\\ chave de upload};
        \node[caixapetroleo,text width=1.5cm,right=of d] (g) {Google Play\\ chave do app};
        \node[caixa,right=of g] (u) {Aparelho};
        \draw[seta] (d) -- node[rotulo,above]{AAB} (g);
        \draw[seta] (g) -- node[rotulo,above]{APKs} (u);
      \end{tikzpicture}
      \begin{dica}
        Perdeu a chave de upload? Com o Play App Signing, o Google permite trocá-la. Sem ele, perder a chave impede atualizar o app.
      \end{dica}
    \end{column}
    \begin{column}{0.48\textwidth}
      \begin{conceito}[iOS]
        Todo app precisa ser assinado para rodar num iPhone. Três peças: o \textbf{certificado} (sua identidade), o \textbf{App ID} e o \textbf{perfil de provisionamento}, que liga certificado, app, capacidades e aparelhos.
      \end{conceito}
      \begin{dica}
        O Xcode gerencia tudo isso com \emph{Automatically manage signing}.
      \end{dica}
      \begin{atencao}
        Chaves e certificados são segredos: nunca no Git; guarde cópias num cofre de senhas.
      \end{atencao}
    \end{column}
  \end{columns}
\end{frame}

\begin{frame}[fragile]{Versionamento}
  \begin{columns}[T]
    \begin{column}{0.5\textwidth}
      \begin{definicao}[Versionamento semântico (SemVer)]
        \textbf{MAJOR.MINOR.PATCH}: \textbf{MAJOR} muda com quebras de compatibilidade, \textbf{MINOR} com funcionalidades novas compatíveis e \textbf{PATCH} com correções.
      \end{definicao}
      \centering
      \begin{tikzpicture}[every node/.style={font=\scriptsize}]
        \node[caixapetroleo] (m) {2};
        \node[right=0.5mm of m] (p1) {.};
        \node[caixaambar,right=0.5mm of p1] (n) {4};
        \node[right=0.5mm of n] (p2) {.};
        \node[caixa,right=0.5mm of p2] (c) {1};
        \node[right=0.5mm of c] (p3) {+};
        \node[caixaescura,right=0.5mm of p3] (b) {27};
        \node[rotulo,below=1mm of m,font=\tiny] {major};
        \node[rotulo,below=1mm of n,font=\tiny] {minor};
        \node[rotulo,below=1mm of c,font=\tiny] {patch};
        \node[rotulo,below=1mm of b,font=\tiny,align=center] {número\\do build};
      \end{tikzpicture}
    \end{column}
    \begin{column}{0.48\textwidth}
      \scriptsize
      \begin{tabularx}{\linewidth}{@{}l L L@{}}
        \toprule
        & \cabecalho{Para pessoas} & \cabecalho{Para a loja} \\
        \midrule
        Android & \cod{versionName} & \cod{versionCode} \\
        iOS & \cod{CFBundleShort\-VersionString} & \cod{CFBundleVersion} \\
        Flutter & \cod{version: 2.4.1+27} & (o número após o \cod{+}) \\
        \bottomrule
      \end{tabularx}
      \begin{atencao}[Regra das lojas]
        O número do build \textbf{sempre aumenta}. Não existe ``voltar versão'': para desfazer um problema, publica-se uma versão nova com número maior.
      \end{atencao}
    \end{column}
  \end{columns}
\end{frame}

\begin{frame}{Formatos de pacote e tamanho do app}
  \begin{columns}[T]
    \begin{column}{0.52\textwidth}
      \scriptsize
      \begin{tabularx}{\linewidth}{@{}l L@{}}
        \toprule
        \cabecalho{Formato} & \cabecalho{O que é} \\
        \midrule
        \textbf{APK} & pacote Android instalável, com tudo para todos os aparelhos \\
        \textbf{AAB} & \emph{Android App Bundle}: formato de \textbf{envio} à Play, que gera APKs sob medida para cada aparelho \\
        \textbf{IPA} & pacote iOS enviado à App Store; a Apple também entrega variantes por aparelho (\emph{App Thinning}) \\
        \bottomrule
      \end{tabularx}
      \centering
      \begin{tikzpicture}[node distance=2mm,every node/.style={font=\tiny}]
        \node[caixapetroleo,minimum width=1.6cm] (aab) {AAB};
        \node[caixa,right=8mm of aab,yshift=5mm] (a1) {arm64 + xxhdpi + pt};
        \node[caixa,right=8mm of aab,yshift=-5mm] (a2) {armeabi + hdpi + en};
        \draw[seta] (aab) -- (a1.west); \draw[seta] (aab) -- (a2.west);
      \end{tikzpicture}
    \end{column}
    \begin{column}{0.46\textwidth}
      \begin{conceito}[Por que o tamanho importa]
        Apps grandes demoram a baixar, ocupam espaço e perdem instalações, principalmente em planos de dados limitados.
      \end{conceito}
      \begin{pordentro}[Como os builds encolhem]
        \textbf{\emph{Tree shaking}} e \textbf{\emph{shrinking}} removem código não usado; a \textbf{ofuscação} troca nomes por símbolos curtos; recursos são comprimidos e separados por aparelho.
      \end{pordentro}
      \begin{atencao}
        Código ofuscado gera relatórios de falha ilegíveis. Guarde os \textbf{símbolos de depuração} de cada versão para decodificá-los.
      \end{atencao}
    \end{column}
  \end{columns}
\end{frame}

\begin{frame}{As duas grandes lojas}
  \scriptsize
  \begin{tabularx}{\textwidth}{@{}p{2.6cm} L L@{}}
    \toprule
    & \cabecalho{Google Play} & \cabecalho{App Store} \\
    \midrule
    Conta & taxa única de US\$\,25; pessoal ou organização & Apple Developer Program, US\$\,99 por ano (isenção para instituições de ensino e sem fins lucrativos) \\
    Revisão & automatizada e humana; de horas a dias (contas novas esperam mais) & humana (\emph{App Review}); em geral de um a dois dias \\
    Testes & faixas interna, fechada e aberta & TestFlight interno (até 100) e externo (até 10\,000) \\
    Lançamento gradual & percentual definido por você & \emph{phased release} automático em 7 dias \\
    Requisito anual & \cod{targetSdk} 36 desde 31/08/2026 & Xcode 26 e SDK do iOS 26 desde 28/04/2026 \\
    Comissão & de 15\,\% a 30\,\% sobre vendas digitais, conforme receita, produto e país & idem \\
    \bottomrule
  \end{tabularx}
  \vspace{4pt}
  \begin{dica}
    Regras de comissão e de pagamento externo mudam com leis e decisões judiciais (União Europeia, EUA, Japão, Coreia). Confira antes de definir o modelo de negócio (Módulo 20).
  \end{dica}
\end{frame}

\begin{frame}{Privacidade desde a concepção e LGPD}
  \begin{columns}[T]
    \begin{column}{0.5\textwidth}
      \begin{definicao}[LGPD (Lei 13.709/2018)]
        Regula o tratamento de \textbf{dados pessoais} no Brasil: tudo que identifica ou pode identificar uma pessoa (nome, e-mail, matrícula, localização, identificador do aparelho).
      \end{definicao}
      \begin{conceito}[Princípios que mais afetam um app]
        \textbf{Finalidade} (para quê), \textbf{necessidade} (o mínimo), \textbf{transparência} (explicar), \textbf{segurança} (proteger) e uma \textbf{base legal} para cada uso (consentimento, execução de contrato\ldots).
      \end{conceito}
    \end{column}
    \begin{column}{0.48\textwidth}
      \begin{conceito}[O que as lojas exigem]
        \begin{itemize}
          \item Política de privacidade publicada.
          \item Declaração de dados coletados: \emph{Data safety} (Play) e \emph{App Privacy} (App Store).
          \item Textos de uso para cada permissão.
          \item \textbf{Exclusão de conta} dentro do app, se ele permite criar conta.
        \end{itemize}
      \end{conceito}
      \begin{atencao}
        A declaração deve incluir o que os SDKs de terceiros coletam (analytics, anúncios, Crashlytics). Não é assessoria jurídica: consulte o encarregado de dados (DPO) da instituição.
      \end{atencao}
    \end{column}
  \end{columns}
\end{frame}

\begin{frame}{Gestão de lançamentos}
  \begin{columns}[c]
    \begin{column}{0.52\textwidth}
      \centering
      \begin{tikzpicture}[x=1.0cm,y=0.45cm]
        \draw[cinzamedio,-{Stealth}] (0,0) -- (6.3,0) node[right,font=\tiny]{dias};
        \draw[cinzamedio,-{Stealth}] (0,0) -- (0,5.6) node[above,font=\tiny]{\% dos usuários};
        \draw[azulpetroleo,line width=1.4pt] (0,0.3) -- (1,0.3) -- (1,1) -- (2,1) -- (2,2.5) -- (3.5,2.5) -- (3.5,5) -- (6,5);
        \node[font=\tiny,azulpetroleo] at (0.5,0.9) {1\%};
        \node[font=\tiny,azulpetroleo] at (1.5,1.6) {5\%};
        \node[font=\tiny,azulpetroleo] at (2.75,3.1) {20\%};
        \node[font=\tiny,azulpetroleo] at (4.75,5.6) {100\%};
        \foreach \x in {1,2,3.5} \fill[ambar] (\x,0) circle (2.2pt);
        \node[rotulo,font=\tiny,align=center] at (2.2,-1.6) {\textcolor{ambar}{$\bullet$} portões: falhas e avaliações\\ dentro do limite? então avança};
      \end{tikzpicture}
    \end{column}
    \begin{column}{0.46\textwidth}
      \begin{definicao}[Lançamento gradual]
        Liberar a versão para uma fração crescente dos usuários, avançando só se os indicadores de saúde continuarem bons.
      \end{definicao}
      \begin{conceito}[Ferramentas do processo]
        \begin{itemize}
          \item \textbf{Trem de versões}: lançar em datas fixas (a cada duas semanas, por exemplo).
          \item \textbf{\emph{Feature flags}}: ligar recursos remotamente, sem nova versão.
          \item \textbf{Versão mínima}: obrigar a atualização quando a antiga é insegura.
          \item \textbf{Notas de versão} claras para usuários e revisores.
        \end{itemize}
      \end{conceito}
    \end{column}
  \end{columns}
\end{frame}

\begin{frame}{ASO: o funil da loja}
  \begin{columns}[c]
    \begin{column}{0.46\textwidth}
      \centering
      \begin{tikzpicture}[every node/.style={font=\scriptsize\bfseries,text=white}]
        \fill[azulpetroleo] (-2.6,3.2) -- (2.6,3.2) -- (2.2,2.5) -- (-2.2,2.5) -- cycle;
        \node at (0,2.85) {Impressões na busca};
        \fill[azulprofundo] (-2.15,2.4) -- (2.15,2.4) -- (1.75,1.7) -- (-1.75,1.7) -- cycle;
        \node at (0,2.05) {Visitas à ficha};
        \fill[ambar!90!black] (-1.7,1.6) -- (1.7,1.6) -- (1.3,0.9) -- (-1.3,0.9) -- cycle;
        \node at (0,1.25) {Instalações};
        \fill[ambar] (-1.25,0.8) -- (1.25,0.8) -- (0.85,0.1) -- (-0.85,0.1) -- cycle;
        \node[text=azulpetroleo] at (0,0.45) {Uso ativo};
      \end{tikzpicture}
    \end{column}
    \begin{column}{0.52\textwidth}
      \begin{definicao}[ASO --- \emph{App Store Optimization}]
        Conjunto de práticas para o app ser \textbf{encontrado} na busca da loja e \textbf{escolhido} por quem vê a ficha.
      \end{definicao}
      \scriptsize
      \begin{tabularx}{\linewidth}{@{}l L@{}}
        \toprule
        \cabecalho{Etapa} & \cabecalho{O que influencia} \\
        \midrule
        Ser encontrado & título, subtítulo ou descrição curta, palavras-chave, categoria \\
        Ser escolhido & ícone, capturas de tela, vídeo, nota e avaliações \\
        Ser mantido & qualidade, desempenho, respostas às avaliações \\
        \bottomrule
      \end{tabularx}
      \begin{dica}
        As duas lojas permitem testes A/B da ficha (ícone, capturas, textos).
      \end{dica}
    \end{column}
  \end{columns}
\end{frame}

\begin{frame}{Depois do lançamento: saúde do app}
  \begin{columns}[T]
    \begin{column}{0.5\textwidth}
      \scriptsize
      \begin{tabularx}{\linewidth}{@{}p{2.3cm} L@{}}
        \toprule
        \cabecalho{Indicador} & \cabecalho{O que mede} \\
        \midrule
        Usuários sem falhas & \% de usuários que não tiveram nenhum \emph{crash} no período \\
        Taxa de ANR & \% de sessões em que a tela travou (Android) \\
        Tempo de inicialização & do toque no ícone à primeira tela útil \\
        Retenção D1/D7/D30 & \% de usuários que voltam 1, 7 e 30 dias após instalar \\
        Nota e avaliações & percepção dos usuários; comentários apontam defeitos \\
        \bottomrule
      \end{tabularx}
    \end{column}
    \begin{column}{0.48\textwidth}
      \begin{atencao}[Limites do Android vitals]
        Acima de 1,09\,\% de falhas percebidas ou 0,47\,\% de ANRs percebidos, o app é considerado de ``mau comportamento'' e perde visibilidade na Play.
      \end{atencao}
      \begin{conceito}[Ciclo contínuo]
        Monitorar $\rightarrow$ priorizar $\rightarrow$ corrigir $\rightarrow$ lançar gradualmente $\rightarrow$ monitorar. Falhas e avaliações alimentam o \emph{backlog} da próxima versão.
      \end{conceito}
      \begin{dica}
        Responda às avaliações com educação e objetividade: a resposta é pública e mostra que o app é cuidado.
      \end{dica}
    \end{column}
  \end{columns}
\end{frame}

% =============================================================================
\section{Fundamentos da distribuição de apps}

\begin{frame}{Do código ao aparelho do usuário}
  \centering
  \begin{tikzpicture}[node distance=2.4mm,every node/.style={font=\scriptsize,minimum height=9mm,minimum width=1.62cm}]
    \node[caixa] (c) {Código\\ \tiny versionado};
    \node[caixapetroleo,right=of c] (b) {Build de\\ release};
    \node[caixapetroleo,right=of b] (a) {Assinatura};
    \node[caixaambar,right=of a] (e) {Envio\\ à loja};
    \node[caixaambar,right=of e] (r) {Revisão\\ e políticas};
    \node[caixa,right=of r] (t) {Faixas\\ de teste};
    \node[caixaescura,right=of t] (p) {Produção};
    \foreach \x/\y in {c/b,b/a,a/e,e/r,r/t,t/p} \draw[seta] (\x) -- (\y);
    \draw[setaambar,rounded corners] (p.south) -- ++(0,-6mm) -| node[rotulo,pos=0.25,below]{\tiny monitorar falhas e avaliações $\rightarrow$ nova versão} (c.south);
  \end{tikzpicture}
  \vspace{6pt}
  \begin{columns}[T]
    \begin{column}{0.48\textwidth}
      \scriptsize
      \begin{tabularx}{\linewidth}{@{}l L L@{}}
        \toprule
        & \cabecalho{Debug} & \cabecalho{Release} \\
        \midrule
        Objetivo & desenvolver & distribuir \\
        Código & com verificações e depuração & otimizado, reduzido e ofuscável \\
        Desempenho & mais lento & o real \\
        Assinatura & chave de depuração automática & sua chave de release \\
        \bottomrule
      \end{tabularx}
    \end{column}
    \begin{column}{0.48\textwidth}
      \begin{definicao}[Build de release]
        Versão do app compilada para o usuário final: sem ferramentas de depuração, com código otimizado e assinada com a chave do desenvolvedor.
      \end{definicao}
      \begin{atencao}
        Teste sempre o build de release antes de publicar: ofuscação e otimizações podem expor defeitos que não aparecem em debug.
      \end{atencao}
    \end{column}
  \end{columns}
\end{frame}

\begin{frame}{Assinatura digital: a teoria}
  \begin{columns}[c]
    \begin{column}{0.54\textwidth}
      \centering
      \begin{tikzpicture}[node distance=4mm,every node/.style={font=\scriptsize}]
        \node[caixa,minimum width=1.6cm] (app) {App};
        \node[caixa,right=of app,minimum width=1.6cm] (h) {Hash\\ \tiny resumo único};
        \node[caixaescura,right=of h,minimum width=2cm] (s) {Assinatura};
        \node[caixaambar,above=5mm of s,minimum width=2cm] (pk) {Chave \textbf{privada}\\ \tiny só o dev tem};
        \draw[seta] (app) -- (h); \draw[seta] (h) -- (s); \draw[seta] (pk) -- (s);
        \node[caixa,below=9mm of app,minimum width=1.6cm] (app2) {App\\ \tiny recebido};
        \node[caixa,right=of app2,minimum width=1.6cm] (h2) {Hash\\ \tiny recalculado};
        \node[caixapetroleo,right=of h2,minimum width=2cm] (v) {Confere?};
        \node[caixa,below=5mm of v,minimum width=2cm] (pub) {Chave \textbf{pública}\\ \tiny no certificado};
        \draw[seta] (app2) -- (h2); \draw[seta] (h2) -- (v); \draw[seta] (pub) -- (v);
        \draw[setaambar] (s.south) -- node[rotulo,right]{\tiny vai junto} (v.north);
        \node[rotulo,left=1mm of app,font=\tiny,rotate=90,anchor=south] {quem publica};
        \node[rotulo,left=1mm of app2,font=\tiny,rotate=90,anchor=south] {sistema};
      \end{tikzpicture}
    \end{column}
    \begin{column}{0.44\textwidth}
      \begin{definicao}[Assinatura digital]
        Resumo (\emph{hash}) do pacote cifrado com a \textbf{chave privada} do desenvolvedor. Qualquer um verifica com a \textbf{chave pública}, mas só quem tem a privada consegue assinar.
      \end{definicao}
      \begin{conceito}[O que ela garante]
        \textbf{Integridade}: o pacote não foi alterado.\\
        \textbf{Autenticidade}: a atualização vem do mesmo autor do app instalado.
      \end{conceito}
      \begin{analogia}
        O lacre de cera com o sinete do remetente: qualquer um reconhece o selo, ninguém o falsifica sem o anel.
      \end{analogia}
    \end{column}
  \end{columns}
\end{frame}

\begin{frame}{Assinatura em cada plataforma}
  \begin{columns}[T]
    \begin{column}{0.5\textwidth}
      \begin{conceito}[Android]
        O pacote é assinado com uma chave guardada num \emph{keystore}. Com o \textbf{Play App Signing}, você assina com a \textbf{chave de upload}; o Google confere e reassina com a \textbf{chave do app}, que fica guardada com ele.
      \end{conceito}
      \centering
      \begin{tikzpicture}[node distance=9mm,every node/.style={font=\tiny,align=center}]
        \node[caixa,text width=1.5cm] (d) {Dev\\ chave de upload};
        \node[caixapetroleo,text width=1.5cm,right=of d] (g) {Google Play\\ chave do app};
        \node[caixa,right=of g] (u) {Aparelho};
        \draw[seta] (d) -- node[rotulo,above]{AAB} (g);
        \draw[seta] (g) -- node[rotulo,above]{APKs} (u);
      \end{tikzpicture}
      \begin{dica}
        Perdeu a chave de upload? Com o Play App Signing, o Google permite trocá-la. Sem ele, perder a chave impede atualizar o app.
      \end{dica}
    \end{column}
    \begin{column}{0.48\textwidth}
      \begin{conceito}[iOS]
        Todo app precisa ser assinado para rodar num iPhone. Três peças: o \textbf{certificado} (sua identidade), o \textbf{App ID} e o \textbf{perfil de provisionamento}, que liga certificado, app, capacidades e aparelhos.
      \end{conceito}
      \begin{dica}
        O Xcode gerencia tudo isso com \emph{Automatically manage signing}.
      \end{dica}
      \begin{atencao}
        Chaves e certificados são segredos: nunca no Git; guarde cópias num cofre de senhas.
      \end{atencao}
    \end{column}
  \end{columns}
\end{frame}

\begin{frame}[fragile]{Versionamento}
  \begin{columns}[T]
    \begin{column}{0.5\textwidth}
      \begin{definicao}[Versionamento semântico (SemVer)]
        \textbf{MAJOR.MINOR.PATCH}: \textbf{MAJOR} muda com quebras de compatibilidade, \textbf{MINOR} com funcionalidades novas compatíveis e \textbf{PATCH} com correções.
      \end{definicao}
      \centering
      \begin{tikzpicture}[every node/.style={font=\scriptsize}]
        \node[caixapetroleo] (m) {2};
        \node[right=0.5mm of m] (p1) {.};
        \node[caixaambar,right=0.5mm of p1] (n) {4};
        \node[right=0.5mm of n] (p2) {.};
        \node[caixa,right=0.5mm of p2] (c) {1};
        \node[right=0.5mm of c] (p3) {+};
        \node[caixaescura,right=0.5mm of p3] (b) {27};
        \node[rotulo,below=1mm of m,font=\tiny] {major};
        \node[rotulo,below=1mm of n,font=\tiny] {minor};
        \node[rotulo,below=1mm of c,font=\tiny] {patch};
        \node[rotulo,below=1mm of b,font=\tiny,align=center] {número\\do build};
      \end{tikzpicture}
    \end{column}
    \begin{column}{0.48\textwidth}
      \scriptsize
      \begin{tabularx}{\linewidth}{@{}l L L@{}}
        \toprule
        & \cabecalho{Para pessoas} & \cabecalho{Para a loja} \\
        \midrule
        Android & \cod{versionName} & \cod{versionCode} \\
        iOS & \cod{CFBundleShort\-VersionString} & \cod{CFBundleVersion} \\
        Flutter & \cod{version: 2.4.1+27} & (o número após o \cod{+}) \\
        \bottomrule
      \end{tabularx}
      \begin{atencao}[Regra das lojas]
        O número do build \textbf{sempre aumenta}. Não existe ``voltar versão'': para desfazer um problema, publica-se uma versão nova com número maior.
      \end{atencao}
    \end{column}
  \end{columns}
\end{frame}

\begin{frame}{Formatos de pacote e tamanho do app}
  \begin{columns}[T]
    \begin{column}{0.52\textwidth}
      \scriptsize
      \begin{tabularx}{\linewidth}{@{}l L@{}}
        \toprule
        \cabecalho{Formato} & \cabecalho{O que é} \\
        \midrule
        \textbf{APK} & pacote Android instalável, com tudo para todos os aparelhos \\
        \textbf{AAB} & \emph{Android App Bundle}: formato de \textbf{envio} à Play, que gera APKs sob medida para cada aparelho \\
        \textbf{IPA} & pacote iOS enviado à App Store; a Apple também entrega variantes por aparelho (\emph{App Thinning}) \\
        \bottomrule
      \end{tabularx}
      \centering
      \begin{tikzpicture}[node distance=2mm,every node/.style={font=\tiny}]
        \node[caixapetroleo,minimum width=1.6cm] (aab) {AAB};
        \node[caixa,right=8mm of aab,yshift=5mm] (a1) {arm64 + xxhdpi + pt};
        \node[caixa,right=8mm of aab,yshift=-5mm] (a2) {armeabi + hdpi + en};
        \draw[seta] (aab) -- (a1.west); \draw[seta] (aab) -- (a2.west);
      \end{tikzpicture}
    \end{column}
    \begin{column}{0.46\textwidth}
      \begin{conceito}[Por que o tamanho importa]
        Apps grandes demoram a baixar, ocupam espaço e perdem instalações, principalmente em planos de dados limitados.
      \end{conceito}
      \begin{pordentro}[Como os builds encolhem]
        \textbf{\emph{Tree shaking}} e \textbf{\emph{shrinking}} removem código não usado; a \textbf{ofuscação} troca nomes por símbolos curtos; recursos são comprimidos e separados por aparelho.
      \end{pordentro}
      \begin{atencao}
        Código ofuscado gera relatórios de falha ilegíveis. Guarde os \textbf{símbolos de depuração} de cada versão para decodificá-los.
      \end{atencao}
    \end{column}
  \end{columns}
\end{frame}

\begin{frame}{As duas grandes lojas}
  \scriptsize
  \begin{tabularx}{\textwidth}{@{}p{2.6cm} L L@{}}
    \toprule
    & \cabecalho{Google Play} & \cabecalho{App Store} \\
    \midrule
    Conta & taxa única de US\$\,25; pessoal ou organização & Apple Developer Program, US\$\,99 por ano (isenção para instituições de ensino e sem fins lucrativos) \\
    Revisão & automatizada e humana; de horas a dias (contas novas esperam mais) & humana (\emph{App Review}); em geral de um a dois dias \\
    Testes & faixas interna, fechada e aberta & TestFlight interno (até 100) e externo (até 10\,000) \\
    Lançamento gradual & percentual definido por você & \emph{phased release} automático em 7 dias \\
    Requisito anual & \cod{targetSdk} 36 desde 31/08/2026 & Xcode 26 e SDK do iOS 26 desde 28/04/2026 \\
    Comissão & de 15\,\% a 30\,\% sobre vendas digitais, conforme receita, produto e país & idem \\
    \bottomrule
  \end{tabularx}
  \vspace{4pt}
  \begin{dica}
    Regras de comissão e de pagamento externo mudam com leis e decisões judiciais (União Europeia, EUA, Japão, Coreia). Confira antes de definir o modelo de negócio (Módulo 20).
  \end{dica}
\end{frame}

\begin{frame}{Privacidade desde a concepção e LGPD}
  \begin{columns}[T]
    \begin{column}{0.5\textwidth}
      \begin{definicao}[LGPD (Lei 13.709/2018)]
        Regula o tratamento de \textbf{dados pessoais} no Brasil: tudo que identifica ou pode identificar uma pessoa (nome, e-mail, matrícula, localização, identificador do aparelho).
      \end{definicao}
      \begin{conceito}[Princípios que mais afetam um app]
        \textbf{Finalidade} (para quê), \textbf{necessidade} (o mínimo), \textbf{transparência} (explicar), \textbf{segurança} (proteger) e uma \textbf{base legal} para cada uso (consentimento, execução de contrato\ldots).
      \end{conceito}
    \end{column}
    \begin{column}{0.48\textwidth}
      \begin{conceito}[O que as lojas exigem]
        \begin{itemize}
          \item Política de privacidade publicada.
          \item Declaração de dados coletados: \emph{Data safety} (Play) e \emph{App Privacy} (App Store).
          \item Textos de uso para cada permissão.
          \item \textbf{Exclusão de conta} dentro do app, se ele permite criar conta.
        \end{itemize}
      \end{conceito}
      \begin{atencao}
        A declaração deve incluir o que os SDKs de terceiros coletam (analytics, anúncios, Crashlytics). Não é assessoria jurídica: consulte o encarregado de dados (DPO) da instituição.
      \end{atencao}
    \end{column}
  \end{columns}
\end{frame}

\begin{frame}{Gestão de lançamentos}
  \begin{columns}[c]
    \begin{column}{0.52\textwidth}
      \centering
      \begin{tikzpicture}[x=1.0cm,y=0.45cm]
        \draw[cinzamedio,-{Stealth}] (0,0) -- (6.3,0) node[right,font=\tiny]{dias};
        \draw[cinzamedio,-{Stealth}] (0,0) -- (0,5.6) node[above,font=\tiny]{\% dos usuários};
        \draw[azulpetroleo,line width=1.4pt] (0,0.3) -- (1,0.3) -- (1,1) -- (2,1) -- (2,2.5) -- (3.5,2.5) -- (3.5,5) -- (6,5);
        \node[font=\tiny,azulpetroleo] at (0.5,0.9) {1\%};
        \node[font=\tiny,azulpetroleo] at (1.5,1.6) {5\%};
        \node[font=\tiny,azulpetroleo] at (2.75,3.1) {20\%};
        \node[font=\tiny,azulpetroleo] at (4.75,5.6) {100\%};
        \foreach \x in {1,2,3.5} \fill[ambar] (\x,0) circle (2.2pt);
        \node[rotulo,font=\tiny,align=center] at (2.2,-1.6) {\textcolor{ambar}{$\bullet$} portões: falhas e avaliações\\ dentro do limite? então avança};
      \end{tikzpicture}
    \end{column}
    \begin{column}{0.46\textwidth}
      \begin{definicao}[Lançamento gradual]
        Liberar a versão para uma fração crescente dos usuários, avançando só se os indicadores de saúde continuarem bons.
      \end{definicao}
      \begin{conceito}[Ferramentas do processo]
        \begin{itemize}
          \item \textbf{Trem de versões}: lançar em datas fixas (a cada duas semanas, por exemplo).
          \item \textbf{\emph{Feature flags}}: ligar recursos remotamente, sem nova versão.
          \item \textbf{Versão mínima}: obrigar a atualização quando a antiga é insegura.
          \item \textbf{Notas de versão} claras para usuários e revisores.
        \end{itemize}
      \end{conceito}
    \end{column}
  \end{columns}
\end{frame}

\begin{frame}{ASO: o funil da loja}
  \begin{columns}[c]
    \begin{column}{0.46\textwidth}
      \centering
      \begin{tikzpicture}[every node/.style={font=\scriptsize\bfseries,text=white}]
        \fill[azulpetroleo] (-2.6,3.2) -- (2.6,3.2) -- (2.2,2.5) -- (-2.2,2.5) -- cycle;
        \node at (0,2.85) {Impressões na busca};
        \fill[azulprofundo] (-2.15,2.4) -- (2.15,2.4) -- (1.75,1.7) -- (-1.75,1.7) -- cycle;
        \node at (0,2.05) {Visitas à ficha};
        \fill[ambar!90!black] (-1.7,1.6) -- (1.7,1.6) -- (1.3,0.9) -- (-1.3,0.9) -- cycle;
        \node at (0,1.25) {Instalações};
        \fill[ambar] (-1.25,0.8) -- (1.25,0.8) -- (0.85,0.1) -- (-0.85,0.1) -- cycle;
        \node[text=azulpetroleo] at (0,0.45) {Uso ativo};
      \end{tikzpicture}
    \end{column}
    \begin{column}{0.52\textwidth}
      \begin{definicao}[ASO --- \emph{App Store Optimization}]
        Conjunto de práticas para o app ser \textbf{encontrado} na busca da loja e \textbf{escolhido} por quem vê a ficha.
      \end{definicao}
      \scriptsize
      \begin{tabularx}{\linewidth}{@{}l L@{}}
        \toprule
        \cabecalho{Etapa} & \cabecalho{O que influencia} \\
        \midrule
        Ser encontrado & título, subtítulo ou descrição curta, palavras-chave, categoria \\
        Ser escolhido & ícone, capturas de tela, vídeo, nota e avaliações \\
        Ser mantido & qualidade, desempenho, respostas às avaliações \\
        \bottomrule
      \end{tabularx}
      \begin{dica}
        As duas lojas permitem testes A/B da ficha (ícone, capturas, textos).
      \end{dica}
    \end{column}
  \end{columns}
\end{frame}

\begin{frame}{Depois do lançamento: saúde do app}
  \begin{columns}[T]
    \begin{column}{0.5\textwidth}
      \scriptsize
      \begin{tabularx}{\linewidth}{@{}p{2.3cm} L@{}}
        \toprule
        \cabecalho{Indicador} & \cabecalho{O que mede} \\
        \midrule
        Usuários sem falhas & \% de usuários que não tiveram nenhum \emph{crash} no período \\
        Taxa de ANR & \% de sessões em que a tela travou (Android) \\
        Tempo de inicialização & do toque no ícone à primeira tela útil \\
        Retenção D1/D7/D30 & \% de usuários que voltam 1, 7 e 30 dias após instalar \\
        Nota e avaliações & percepção dos usuários; comentários apontam defeitos \\
        \bottomrule
      \end{tabularx}
    \end{column}
    \begin{column}{0.48\textwidth}
      \begin{atencao}[Limites do Android vitals]
        Acima de 1,09\,\% de falhas percebidas ou 0,47\,\% de ANRs percebidos, o app é considerado de ``mau comportamento'' e perde visibilidade na Play.
      \end{atencao}
      \begin{conceito}[Ciclo contínuo]
        Monitorar $\rightarrow$ priorizar $\rightarrow$ corrigir $\rightarrow$ lançar gradualmente $\rightarrow$ monitorar. Falhas e avaliações alimentam o \emph{backlog} da próxima versão.
      \end{conceito}
      \begin{dica}
        Responda às avaliações com educação e objetividade: a resposta é pública e mostra que o app é cuidado.
      \end{dica}
    \end{column}
  \end{columns}
\end{frame}


% =============================================================================
\section{Preparando o app Flutter}

\begin{frame}[fragile]{Identidade: nome, ícone, \emph{splash} e versão}
  \begin{columns}[T]
    \begin{column}{0.54\textwidth}
\begin{codigo}[title=pubspec.yaml]{YAML}
name: labmonitor
version: 1.2.0+12        # versão + número do build

flutter_launcher_icons:
  android: true
  ios: true
  image_path: assets/icone/labmonitor.png
  adaptive_icon_background: "#0D3052"
  adaptive_icon_foreground: assets/icone/frente.png

flutter_native_splash:
  color: "#0D3052"
  image: assets/icone/splash.png
\end{codigo}
\begin{terminal}
$ flutter pub add -d flutter_launcher_icons \
    flutter_native_splash
$ dart run flutter_launcher_icons
$ dart run flutter_native_splash:create
\end{terminal}
    \end{column}
    \begin{column}{0.44\textwidth}
      \scriptsize
      \begin{tabularx}{\linewidth}{@{}l L@{}}
        \toprule
        \cabecalho{Item} & \cabecalho{Onde fica} \\
        \midrule
        ID Android & \cod{applicationId} em \cod{android/app/build.gradle.kts} \\
        ID iOS & \emph{Bundle Identifier} no alvo \cod{Runner} do Xcode \\
        Nome exibido & \cod{android:label} (manifesto) e \cod{CFBundleDisplayName} (\cod{Info.plist}) \\
        Versão & \cod{version} do \cod{pubspec.yaml} \\
        \bottomrule
      \end{tabularx}
      \begin{atencao}
        O identificador (\cod{br.edu.ifma.labmonitor}) é \textbf{permanente} depois da primeira publicação. Escolha com cuidado.
      \end{atencao}
    \end{column}
  \end{columns}
\end{frame}

\begin{frame}[fragile]{Modos de compilação e builds de release}
  \begin{columns}[T]
    \begin{column}{0.46\textwidth}
      \scriptsize
      \begin{tabularx}{\linewidth}{@{}l L@{}}
        \toprule
        \cabecalho{Modo} & \cabecalho{Características} \\
        \midrule
        \textbf{debug} & compilação JIT, \emph{hot reload}, \cod{assert}s ativos, DevTools completo \\
        \textbf{profile} & AOT como o release, com medição de desempenho \\
        \textbf{release} & AOT, \emph{tree shaking}, sem \cod{assert}s nem depuração \\
        \bottomrule
      \end{tabularx}
      \begin{pordentro}[Por que o release é mais rápido]
        Em release, o Dart é compilado \textbf{antes} (AOT) para código de máquina ARM; em debug, roda numa VM com compilação sob demanda (JIT), que permite o \emph{hot reload}.
      \end{pordentro}
    \end{column}
    \begin{column}{0.52\textwidth}
\begin{terminal}[]
# Android: App Bundle para a Play
$ flutter build appbundle --release \
    --obfuscate --split-debug-info=build/simbolos

# iOS: pacote para a App Store (no Mac)
$ flutter build ipa --release \
    --obfuscate --split-debug-info=build/simbolos

# o que ocupa espaço no app?
$ flutter build appbundle --analyze-size
\end{terminal}
      \begin{atencao}
        Guarde a pasta \cod{build/simbolos} de \textbf{cada} versão e envie-a ao Crashlytics (\cod{firebase crashlytics:symbols:upload}); sem ela, as pilhas de erro ficam ilegíveis.
      \end{atencao}
    \end{column}
  \end{columns}
\end{frame}

\begin{frame}[fragile]{Ambientes: desenvolvimento, homologação e produção}
  \begin{columns}[T]
    \begin{column}{0.54\textwidth}
\begin{codigo}[title=env/producao.json]{JSON}
{
  "API_URL": "https://api.labmonitor.dev/v1",
  "AMBIENTE": "producao"
}
\end{codigo}
\begin{codigo}[title=lib/config.dart]{Dart}
const apiUrl = String.fromEnvironment('API_URL',
    defaultValue: 'http://10.0.2.2:8080/v1');
const ambiente = String.fromEnvironment('AMBIENTE',
    defaultValue: 'dev');
\end{codigo}
\begin{terminal}
$ flutter run --dart-define-from-file=env/dev.json
$ flutter build appbundle \
    --dart-define-from-file=env/producao.json
\end{terminal}
    \end{column}
    \begin{column}{0.44\textwidth}
      \begin{conceito}[Mesmo código, configurações diferentes]
        A URL da API, o projeto Firebase e o nível de log mudam por ambiente. O valor é decidido na hora do build, não no código.
      \end{conceito}
      \begin{dica}[\emph{Flavors}]
        Para ter ``LabMonitor Homolog'' e ``LabMonitor'' instalados lado a lado, use \emph{flavors} (\cod{productFlavors} no Android, \emph{schemes} no iOS) e \cod{--flavor homolog}.
      \end{dica}
      \begin{erro}[Segredo no binário]
        Valores de \cod{--dart-define} ficam gravados no app e podem ser extraídos.
      \end{erro}
    \end{column}
  \end{columns}
\end{frame}

% =============================================================================
\section{Google Play}

\begin{frame}[fragile]{Chave de upload e assinatura no Gradle}
  \begin{columns}[T]
    \begin{column}{0.56\textwidth}
\begin{codigo}[title=android/app/build.gradle.kts (trechos)]{Kotlin}
val props = Properties().apply {
    val f = rootProject.file("key.properties")
    if (f.exists()) load(FileInputStream(f))
}

android {
    signingConfigs {
        create("release") {
            keyAlias = props["keyAlias"] as String
            keyPassword = props["keyPassword"] as String
            storeFile = file(props["storeFile"] as String)
            storePassword = props["storePassword"] as String
        }
    }
    buildTypes {
        release {
            signingConfig = signingConfigs.getByName("release")
        }
    }
}
\end{codigo}
    \end{column}
    \begin{column}{0.42\textwidth}
\begin{terminal}
$ keytool -genkey -v -keystore \
    ~/chaves/upload.jks \
    -keyalg RSA -keysize 2048 \
    -validity 10000 -alias upload
\end{terminal}
\begin{codigo}[title=android/key.properties]{bash}
storePassword=********
keyPassword=********
keyAlias=upload
storeFile=/home/aluno/chaves/upload.jks
\end{codigo}
      \begin{atencao}
        \cod{key.properties} e \cod{*.jks} vão para o \cod{.gitignore}. No CI, use segredos do repositório.
      \end{atencao}
    \end{column}
  \end{columns}
\end{frame}

\begin{frame}{Play Console: do cadastro à produção}
  \begin{columns}[T]
    \begin{column}{0.5\textwidth}
      \begin{conceito}[Antes do primeiro envio]
        \begin{enumerate}
          \item Criar o app (nome, idioma, gratuito ou pago).
          \item Ficha da loja: título (até 30 caracteres), descrição curta (80), descrição completa, ícone 512$\times$512, destaque 1024$\times$500, capturas.
          \item Classificação indicativa (questionário IARC).
          \item Segurança dos dados, público-alvo, anúncios.
          \item Acesso ao app: credenciais de teste para os revisores.
          \item Enviar o AAB às faixas: \textbf{interna} (até 100 pessoas) $\rightarrow$ \textbf{fechada} $\rightarrow$ \textbf{aberta} $\rightarrow$ \textbf{produção}.
        \end{enumerate}
      \end{conceito}
    \end{column}
    \begin{column}{0.48\textwidth}
      \begin{atencao}[Regras atuais]
        \begin{itemize}
          \item \cod{targetSdk} 36 ou maior desde 31/08/2026 (confira em \cod{build.gradle.kts}).
          \item Contas \textbf{pessoais} criadas após 13/11/2023: teste fechado com 12 testadores por 14 dias seguidos.
          \item Desde 30/09/2026, no Brasil, apps em aparelhos certificados precisam vir de desenvolvedores verificados, mesmo fora da Play.
        \end{itemize}
      \end{atencao}
    \end{column}
  \end{columns}
\end{frame}

% =============================================================================
\section{App Store}

\begin{frame}{Pré-requisitos da Apple}
  \begin{columns}[T]
    \begin{column}{0.5\textwidth}
      \begin{conceito}[O que você precisa]
        \begin{itemize}
          \item Um \textbf{Mac} com \textbf{Xcode 26} ou mais novo (exigido para envios desde 28/04/2026).
          \item Conta no \textbf{Apple Developer Program}.
          \item \emph{Bundle ID} registrado e assinatura automática no alvo \cod{Runner} (\emph{Signing \& Capabilities}).
          \item Registro do app no \textbf{App Store Connect}: nome, idioma, \emph{bundle ID} e SKU.
        \end{itemize}
      \end{conceito}
      \begin{dica}[Sem Mac?]
        Serviços de CI com máquinas macOS (Codemagic, GitHub Actions, Bitrise) geram e enviam o build iOS.
      \end{dica}
    \end{column}
    \begin{column}{0.48\textwidth}
      \begin{conceito}[Privacidade no iOS]
        \begin{itemize}
          \item Textos de uso no \cod{Info.plist} para cada permissão (Módulo 13).
          \item \emph{App Privacy} preenchido no App Store Connect.
          \item \textbf{Manifesto de privacidade} (\cod{PrivacyInfo.xcprivacy}) declarando dados coletados e o motivo do uso de certas APIs; plugins trazem os seus.
        \end{itemize}
      \end{conceito}
      \begin{dica}[Instituições]
        Instituições de ensino podem solicitar isenção da anuidade do programa.
      \end{dica}
    \end{column}
  \end{columns}
\end{frame}

\begin{frame}[fragile]{Build, envio e TestFlight}
  \begin{columns}[T]
    \begin{column}{0.52\textwidth}
\begin{terminal}
$ flutter build ipa --release \
    --obfuscate --split-debug-info=build/simbolos
# gera build/ios/ipa/*.ipa

# envio pelo terminal (chave da API do
# App Store Connect) ou pelo app Transporter
$ xcrun altool --upload-app --type ios \
    -f build/ios/ipa/*.ipa \
    --apiKey $CHAVE --apiIssuer $EMISSOR
\end{terminal}
      \begin{dica}
        Alternativa visual: abrir \cod{ios/Runner.xcworkspace} no Xcode, \emph{Product} $\rightarrow$ \emph{Archive} $\rightarrow$ \emph{Distribute App}.
      \end{dica}
    \end{column}
    \begin{column}{0.46\textwidth}
      \scriptsize
      \begin{tabularx}{\linewidth}{@{}l L@{}}
        \toprule
        \cabecalho{TestFlight} & \cabecalho{Como funciona} \\
        \midrule
        Interno & até 100 membros da equipe; sem revisão \\
        Externo & até 10\,000 pessoas por e-mail ou link público; passa por revisão beta \\
        Validade & cada build expira em 90 dias \\
        \bottomrule
      \end{tabularx}
      \begin{conceito}[Envio para revisão]
        Escolha o build, preencha as informações para o revisor (conta de demonstração!) e envie. Aprovado, publique na hora, numa data ou com \emph{phased release}.
      \end{conceito}
    \end{column}
  \end{columns}
\end{frame}

\begin{frame}{App Review: motivos comuns de rejeição}
  \footnotesize
  \begin{tabularx}{\textwidth}{@{}p{3.6cm} L L@{}}
    \toprule
    \cabecalho{Diretriz} & \cabecalho{Problema típico} & \cabecalho{Como evitar} \\
    \midrule
    2.1 Completude & falhas, telas vazias, \emph{links} quebrados, conteúdo provisório & testar o build de release em aparelho real \\
    2.3 Metadados & capturas que não mostram o app real, descrição enganosa & imagens fiéis às telas atuais \\
    4.2 Funcionalidade mínima & app que só embrulha um site & oferecer valor de app nativo \\
    4.8 Serviços de login & login social sem opção equivalente focada em privacidade & incluir Sign in with Apple \\
    5.1.1 Coleta de dados & permissão sem texto claro; conta sem exclusão & textos específicos; excluir conta no app \\
    3.1.1 Compras no app & vender conteúdo digital fora do sistema da Apple & usar compra no app onde for exigido \\
    \bottomrule
  \end{tabularx}
  \vspace{4pt}
  \begin{dica}
    Se o app exige login, informe uma conta de demonstração. Responda à revisão pelo \emph{Resolution Center}, com calma e objetividade.
  \end{dica}
\end{frame}

% =============================================================================
\section{Automação e operação}

\begin{frame}[fragile]{Automatizando a entrega com fastlane}
  \begin{columns}[T]
    \begin{column}{0.58\textwidth}
\begin{codigo}[title=android/fastlane/Fastfile e ios/fastlane/Fastfile]{bash}
# Android
lane :interno do
  sh("cd ../.. && flutter build appbundle --release")
  upload_to_play_store(track: "internal",
    aab: "../build/app/outputs/bundle/release/app-release.aab")
end

# iOS
lane :beta do
  sh("cd ../.. && flutter build ipa --release")
  upload_to_testflight(
    ipa: "../build/ios/ipa/LabMonitor.ipa")  # nome do app
end
\end{codigo}
\begin{terminal}
$ cd android && fastlane interno
$ cd ios && fastlane beta
\end{terminal}
    \end{column}
    \begin{column}{0.4\textwidth}
      \begin{conceito}[Entrega contínua]
        Cada versão aprovada no CI é compilada, assinada e enviada à faixa de teste sem passos manuais. Menos erro humano, entregas frequentes.
      \end{conceito}
      \begin{dica}
        O fastlane também gera capturas de tela, atualiza a ficha da loja e incrementa o número do build.
      \end{dica}
      \begin{dica}
        O Módulo 18 monta o \emph{pipeline} completo no GitHub Actions.
      \end{dica}
    \end{column}
  \end{columns}
\end{frame}

\begin{frame}[fragile]{Atualização obrigatória e pedido de avaliação}
  \begin{columns}[T]
    \begin{column}{0.58\textwidth}
\begin{codigo}[title=Na abertura do app]{Dart}
// versão mínima no Remote Config (Módulo 12)
final minima = Version.parse(rc.getString('versao_minima'));
final info = await PackageInfo.fromPlatform();
if (Version.parse(info.version) < minima) {
  router.go('/atualizar');          // bloqueia o uso
}

// Android: atualização pela própria Play
final upd = await InAppUpdate.checkForUpdate();
if (upd.updateAvailability ==
    UpdateAvailability.updateAvailable) {
  await InAppUpdate.performImmediateUpdate();
}

// após uma reserva bem-sucedida
final review = InAppReview.instance;
if (await review.isAvailable()) review.requestReview();
\end{codigo}
    \end{column}
    \begin{column}{0.4\textwidth}
      \begin{conceito}[Pacotes]
        \cod{pub\_semver} compara versões; \cod{in\_app\_update} usa a API de atualização da Play (só Android); \cod{in\_app\_review} abre a avaliação nativa das duas lojas.
      \end{conceito}
      \begin{atencao}
        O sistema decide se mostra o pedido de avaliação (no iOS, no máximo três vezes por ano). Nunca ofereça recompensa por nota nem peça logo na primeira abertura.
      \end{atencao}
    \end{column}
  \end{columns}
\end{frame}

% =============================================================================
\section{Fechamento}

\begin{frame}{Resumo}
  \begin{resumo}
    \begin{itemize}
      \item O release é otimizado, ofuscado e assinado; a assinatura garante integridade e autoria.
      \item Versão para pessoas (SemVer) e número de build para a loja, sempre crescente.
      \item \cod{flutter build appbundle} e \cod{flutter build ipa} geram os pacotes; guarde os símbolos.
      \item Play: faixas de teste, regras de \cod{targetSdk}, testadores e verificação de desenvolvedor.
      \item App Store: Xcode 26, TestFlight, App Review e declarações de privacidade.
      \item Lançamento gradual, monitoramento e automação fecham o ciclo.
    \end{itemize}
  \end{resumo}
\end{frame}

\begin{frame}{Glossário (1/2)}
  \begin{glossario}
    Build de release & versão otimizada e assinada do app, destinada ao usuário final \\
    Assinatura digital & \emph{hash} do pacote cifrado com a chave privada; prova integridade e autoria \\
    \emph{Keystore} & arquivo protegido por senha que guarda as chaves de assinatura no Android \\
    Play App Signing & serviço em que o Google guarda a chave do app e você assina com a chave de upload \\
    Perfil de provisionamento & arquivo da Apple que liga certificado, App ID, capacidades e aparelhos \\
    AAB & \emph{Android App Bundle}: formato de envio do qual a Play gera APKs por aparelho \\
    IPA & pacote do app iOS enviado ao App Store Connect \\
    Ofuscação & troca de nomes do código por símbolos curtos, dificultando engenharia reversa \\
  \end{glossario}
\end{frame}

\begin{frame}{Glossário (2/2)}
  \begin{glossario}
    Símbolos de depuração & arquivos que traduzem pilhas ofuscadas de volta para nomes legíveis \\
    SemVer & versionamento semântico: MAJOR.MINOR.PATCH \\
    Faixa de teste & canal da Play para distribuir a grupos (interna, fechada, aberta) \\
    TestFlight & serviço da Apple para distribuir builds de teste \\
    App Review & revisão humana dos apps enviados à App Store \\
    Lançamento gradual & liberação da versão para uma porcentagem crescente de usuários \\
    ASO & \emph{App Store Optimization}: práticas para ser encontrado e escolhido na loja \\
    \emph{Flavor} & variante do app (homologação, produção) gerada do mesmo código \\
  \end{glossario}
\end{frame}

\begin{frame}{Questões de revisão}
  \begin{questoes}
    \item O que a assinatura digital garante, e por que perder a chave pode impedir atualizações?
    \item Qual a diferença entre a chave de upload e a chave do app no Play App Signing?
    \item Por que o número do build precisa sempre aumentar? Como ``desfazer'' uma versão com defeito?
    \item Em que o AAB difere do APK, e que vantagem ele traz para o usuário?
    \item Por que é preciso guardar os símbolos de depuração de cada versão ofuscada?
    \item Por que valores passados com \cod{--dart-define} não servem para guardar segredos?
    \item Cite três motivos comuns de rejeição na App Store e como evitá-los.
    \item Como um lançamento gradual reduz o risco de uma versão nova?
  \end{questoes}
\end{frame}

\begin{frame}{Exercícios}
  \begin{exercicio}[Exercício 1 --- Identidade]
    Defina o identificador, gere ícone adaptativo e \emph{splash} com as cores do curso e configure três ambientes com \cod{--dart-define-from-file}.
  \end{exercicio}
  \begin{exercicio}[Exercício 2 --- Release Android]
    Crie a chave de upload, configure a assinatura, gere o AAB ofuscado e analise o tamanho com \cod{--analyze-size}.
  \end{exercicio}
  \begin{exercicio}[Exercício 3 --- Ficha da loja]
    Escreva título, descrição curta e completa e planeje cinco capturas de tela, aplicando as ideias de ASO.
  \end{exercicio}
  \begin{projeto}[LabMonitor --- Etapa 15]
    Publique o LabMonitor na faixa de teste fechado da Play com a turma (e no TestFlight, se houver Mac), com política de privacidade, exclusão de conta e lançamento gradual da versão 1.0.
  \end{projeto}
\end{frame}

\begin{frame}{Para aprofundar}
  \small
  \begin{itemize}
    \item Publicar no Android: \url{https://docs.flutter.dev/deployment/android}
    \item Publicar no iOS: \url{https://docs.flutter.dev/deployment/ios}
    \item Ofuscação: \url{https://docs.flutter.dev/deployment/obfuscate}
    \item Play Console: \url{https://support.google.com/googleplay/android-developer}
    \item Diretrizes da App Store: \url{https://developer.apple.com/app-store/review/guidelines}
    \item fastlane: \url{https://docs.fastlane.tools}
  \end{itemize}
\end{frame}

\begin{frame}{Respostas das questões de revisão}
  \begin{respostas}
    \item Integridade (o pacote não foi alterado) e autenticidade (vem do mesmo autor). O sistema só aceita atualizações assinadas pela mesma chave; sem ela, não há como provar a autoria.
    \item A de upload assina o que você envia e pode ser trocada; a do app, guardada pelo Google, assina os APKs entregues aos aparelhos.
    \item A loja e o sistema usam o número para saber qual build é mais novo. Para desfazer, publica-se uma versão nova, com número maior, contendo a correção (ou o código anterior).
    \item O APK contém tudo para todos os aparelhos; o AAB permite à Play gerar APKs só com o código, as imagens e os idiomas de cada aparelho: download menor.
    \item Sem eles, as pilhas de falha mostram nomes ofuscados e não é possível localizar o defeito no código.
    \item Eles são gravados no binário, que pode ser descompactado e analisado por qualquer pessoa.
    \item Falhas ou telas vazias (testar o release), login social sem opção equivalente (Sign in with Apple), permissões sem texto claro ou conta sem exclusão (textos específicos e exclusão no app).
    \item Um defeito atinge só a fração inicial de usuários; os indicadores de saúde decidem se o lançamento avança ou é interrompido.
  \end{respostas}
\end{frame}

\encerramento{Módulo 16 (comum) --- Backend terceirizado: Firebase, Supabase, AWS Amplify, Appwrite e outros.}

\end{document}
